\section{Anisotropy: thermal coupling of quasi-longitudinal waves in cubic crystals}\label{sec:aniso}

The theory of \cref{sec:hyperbolic} is isotropic; a single crystal is not. Devices use single crystals, so the questions are how the coupling strength $\varepsilon$ depends on propagation direction and which modes it touches.

\subsection{Christoffel problem with thermal coupling}
For a cubic crystal the thermal expansion is isotropic, $\alpha_{ij}=\alpha_T\delta_{ij}$, so the thermal-stress tensor is $\beta_{ij}=\beta_0\delta_{ij}$ with $\beta_0=(c_{11}+2c_{12})\alpha_T$. Heat conduction is scalar. For a plane wave with direction $\bm n$ the field equations give, in the low-frequency (adiabatic) limit $\omega\ll\omega_*$, the Christoffel equation
\begin{equation}\label{eq:christ}
\bigl[\Gamma_{ik}(\bm n)+\tfrac{\beta_0^2T_0}{c_v}\,n_in_k-\rho v^2\delta_{ik}\bigr]p_k=0,\qquad \Gamma_{ik}=C_{ijkl}n_jn_l,
\end{equation}
with $c_v$ the volumetric heat capacity; the high-frequency (isothermal) limit is the same equation without the thermal term. The thermal term is a rank-one matrix along $\bm n$, so it couples only the polarization component along the propagation direction and vanishes for purely transverse modes. Along the symmetry directions of a cubic crystal the coupling is therefore governed by
\begin{equation}\label{eq:epsdir}
\varepsilon_{\bm n}=\frac{\beta_0^2T_0}{c_v\,M_{\bm n}},\qquad M_{[100]}=c_{11},\ M_{[110]}=\tfrac12(c_{11}+c_{12}+2c_{44}),\ M_{[111]}=\tfrac13(c_{11}+2c_{12}+4c_{44}),
\end{equation}
with $M_{\bm n}$ the longitudinal modulus along $\bm n$; the quasi-longitudinal speed rises from the isothermal to the adiabatic value by the factor $\sqrt{1+\varepsilon_{\bm n}}$, and both transverse modes along these directions are untouched.

\subsection{Numerical values}
\Cref{tab:aniso} gives the result for three cubic metals and semiconductors, from the anisotropy figure (\cref{fig:aniso}) with rounded handbook values of the elastic constants, expansion coefficient, density and specific heat at $T_0=300$\,K; a specific sample needs its own measurements. The speed shifts from the full eigenvalue problem agree with $\varepsilon_{\bm n}/2$ to the digits shown.

\begin{table}[ht]\centering\small
\caption{Direction dependence of the thermoelastic coupling of quasi-longitudinal waves ($T_0=300$\,K, rounded handbook inputs). $\varepsilon_{\bm n}$ from \cref{eq:epsdir}; shift is $100(v_{\rm ad}/v_{\rm iso}-1)$.}\label{tab:aniso}
\begin{tabular}{@{}lccccc@{}}\toprule
 & \multicolumn{2}{c}{$[100]$} & \multicolumn{2}{c}{$[111]$} & ratio\\
Crystal & $\varepsilon$ & shift (\%) & $\varepsilon$ & shift (\%) & $\varepsilon_{[100]}/\varepsilon_{[111]}$\\\midrule
Cu & 0.0238 & 1.18 & 0.0169 & 0.84 & 1.41\\
Al & 0.0325 & 1.61 & 0.0306 & 1.52 & 1.06\\
Si & $6.5\times10^{-4}$ & 0.032 & $5.3\times10^{-4}$ & 0.026 & 1.23\\
\bottomrule\end{tabular}
\end{table}

For copper the coupling differs by $41\%$ between $[100]$ and $[111]$. The thermoelastic loss peak, proportional to $\varepsilon_{\bm n}$ in the weak-coupling Debye form, and the adiabatic--isothermal speed gap therefore vary with direction by a comparable amount. Aluminum has only $6\%$ coupling anisotropy, because its elastic anisotropy is weak. Silicon has a tiny coupling at all angles, so longitudinal-wave thermoelastic loss in silicon is very small.

\begin{figure}[t]
\centering
\includegraphics[width=\linewidth]{figs/fig_aniso.pdf}
\caption{(a) Adiabatic phase speeds of the three modes of copper in the (001) plane as functions of the propagation angle from $[100]$. (b) Relative stiffening $v_{\rm ad}/v_{\rm iso}-1$ of the quasi-longitudinal mode for Cu, Si and Al.}
\label{fig:aniso}
\end{figure}

\subsection{What changes for the heat-flux law}
In non-cubic crystals anisotropic conduction makes $G$ a tensor: the template of \cref{sec:foundations} generalizes to $\hat q_i=-G_{ij}(s)\hat\theta_{,j}$, and lag times may differ from axis to axis. For the cubic case used here conduction is isotropic, so only elastic and coupling anisotropy enter, and every scalar result of the preceding sections carries over with $c_1^2\to M_{\bm n}/\rho$ and $\varepsilon\to\varepsilon_{\bm n}$ along a high-symmetry direction. Away from them the quasi-longitudinal polarization is not along $\bm n$, the coupling falls by $(\bm n\!\cdot\!\bm p)^2$, and a quasi-transverse mode picks up a small thermal component.


\begin{summarybox}[Summary]
Anisotropy replaces scalar moduli by tensors $C_{ijkl},\beta_{ij},k_{ij}$; shear and dilatation no longer decouple and the Stroh formalism is the systematic tool. Along high-symmetry directions scalar results carry over with effective moduli. Away from them, quasi-longitudinal and quasi-shear modes mix, coupling weakens by $(\bm n\cdot\bm p)^2$, and quasi-transverse waves acquire a small thermal component.
\end{summarybox}